\section{A total category whose colimits are not exact}\label{sec:cdelta}

\subsection{The category $\Cd$}

A \emph{hypergraph} is a finite set $X$ (of vertices) with a set $E_X$ of
three-element subsets of $X$ (the hyperedges). A subset $U\subseteq X$
carries the induced hypergraph, with hyperedges
$E_X(U)=\{h\in E_X: h\subseteq U\}$. Put
\[
 \delta_X(U)=2\,|U|-3\,|E_X(U)| .
\]
We write $\delta(U)$ when $X$ is understood; $\delta(U)$ depends only on the
induced hypergraph on $U$. The hypergraph $X$ is \emph{positive} if
$\delta(U)>0$ for every non-empty $U\subseteq X$. For $A\subseteq B\subseteq X$
we say that $A$ is \emph{closed in} $B$, and write $A\cle B$, if
\[
 \delta(U)>\delta(A)\qquad\text{for every } U \text{ with } A\subsetneq U\subseteq B .
\]
Thus $X$ is positive if and only if $\varnothing\cle X$. An \emph{ordered}
hypergraph carries in addition a linear order on its vertices. An
\emph{embedding} $f\colon X\to Y$ of ordered hypergraphs is an increasing
injection such that a three-element subset $h$ of $X$ is a hyperedge of $X$
exactly when $f(h)$ is a hyperedge of $Y$; it is a \emph{closed embedding} if
$f(X)\cle Y$.

\begin{definition}\label{def:cdelta}
The objects of $\Cd$ are the positive ordered hypergraphs with vertex set
$[n]=\{1<\dots<n\}$, $n\ge 1$. Its arrows are the closed embeddings, composed as
maps.
\end{definition}

Composites of closed embeddings are closed by Lemma~\ref{lem:calc}(c) below, so
$\Cd$ is a category; it is small, and every non-empty positive ordered
hypergraph is isomorphic to exactly one object. We write $V$ for the object
with one vertex and $W$ for the object $([3],\{\{1,2,3\}\})$. The latter is
positive: its subsets with one, two and three elements have
$\delta=2,4,3$. The same numbers show that the three maps
$v_1,v_2,v_3\colon V\to W$, $v_i(1)=i$, are closed embeddings.

\begin{lemma}[closure calculus]\label{lem:calc}
Let $X$ be a hypergraph.
\begin{enumerate}
\item[\textup{(a)}] $\delta(U)+\delta(U')\ge\delta(U\cup U')+\delta(U\cap U')$ for all
 $U,U'\subseteq X$.
\item[\textup{(b)}] If $A\cle B$ and $Y\subseteq B$, then $A\cap Y\cle Y$.
\item[\textup{(c)}] If $A\cle B$ and $B\cle C$, then $A\cle C$.
\item[\textup{(d)}] Let $X=B\cup C$ and $A=B\cap C$, and suppose that
 $E_X=E_X(B)\cup E_X(C)$ \textup{(}we then call $X$ the \emph{free amalgam} of $B$
 and $C$ over $A$\textup{)}. Then for every $U\subseteq X$
 \begin{equation}\label{eq:amalgam}
  \delta(U)=\delta(U\cap B)+\delta(U\cap C)-\delta(U\cap A).
 \end{equation}
 If moreover $B$ and $C$ are positive, $A\cle B$ and $A\cle C$, then $X$ is
 positive, $B\cle X$ and $C\cle X$. Here $A$ may be empty.
\end{enumerate}
\end{lemma}

\begin{proof}
(a) The number of vertices is modular. For the hyperedges,
$E(U)\cup E(U')\subseteq E(U\cup U')$ and $E(U)\cap E(U')=E(U\cap U')$, so
$|E(U)|+|E(U')|\le|E(U\cup U')|+|E(U\cap U')|$.

(b) Let $A\cap Y\subsetneq U\subseteq Y$. Then $U\cap A=A\cap Y$ and
$U\cup A\supsetneq A$, so by (a)
$\delta(U)-\delta(A\cap Y)\ge\delta(U\cup A)-\delta(A)>0$.

(c) Let $A\subsetneq U\subseteq C$. If $U\subseteq B$ there is nothing to prove.
Otherwise $U\cap B\subsetneq U$ and $U\cap B\cle U$ by (b), so
$\delta(U)>\delta(U\cap B)\ge\delta(A)$, the last inequality because
$A\subseteq U\cap B\subseteq B$ and $A\cle B$.

(d) A hyperedge of $X$ lies in $B$ or in $C$, and those lying in both are the
hyperedges contained in $A$; this gives \eqref{eq:amalgam}. Let $U$ be
non-empty. If $U\subseteq B$ then $\delta(U)>0$. Otherwise
$U\cap A\subsetneq U\cap C$, and $U\cap A\cle U\cap C$ by (b), so
$\delta(U\cap C)-\delta(U\cap A)>0$, while $\delta(U\cap B)\ge 0$; hence
$\delta(U)>0$. If $B\subsetneq U\subseteq X$, then $U\cap B=B$, $U\cap A=A$ and
$U\cap C\supsetneq A$, so
$\delta(U)-\delta(B)=\delta(U\cap C)-\delta(A)>0$. Thus $B\cle X$, and
$C\cle X$ by symmetry.
\end{proof}

If, in (d), $B$ and $C$ are ordered and their orders agree on $A$, the free
amalgam carries linear orders extending both. We call \emph{the ordered free
amalgam with $B$ first} the one in which, for $b\in B\setminus A$ and
$c\in C\setminus A$, we have $b<c$ exactly when every element of $A$ below $b$
is below $c$.

\subsection{The generic structure and its automorphism group}

For an ordered hypergraph $M$ with countably many vertices, all of whose finite
subsets are positive, and a finite $B\subseteq M$, we write $B\cle M$ if
$B\cle U$ for every finite $U$ with $B\subseteq U\subseteq M$. Parts (b) and
(c) of Lemma~\ref{lem:calc} remain true with $M$ in the place of the largest
set.

\begin{proposition}\label{prop:generic}
There is a countable ordered hypergraph $M$, all of whose finite subsets are
positive, with the following properties.
\begin{enumerate}
\item[\textup{(M1)}] Every finite subset of $M$ is contained in a finite subset
 $B$ with $B\cle M$.
\item[\textup{(M2)}] If $B\cle M$ is finite and $f\colon B\to C$ is a closed
 embedding into a finite positive ordered hypergraph, there is an embedding
 $h\colon C\to M$ with $h(C)\cle M$ and $hf=\id_B$.
\item[\textup{(M3)}] Every isomorphism between two finite closed subsets of $M$
 extends to an automorphism of $M$.
\end{enumerate}
\end{proposition}

\begin{proof}
A \emph{task} is a triple $(k,B,f)$ where $k\ge 0$ and $f\colon B\to C$ is a
closed embedding of a subset $B$ of the $k$-th stage, closed in that stage,
into a finite positive ordered hypergraph. Up to isomorphism of $C$ under $B$
there are countably many tasks at each stage, so the tasks of all stages can
be enumerated in such a way that every task is met at a later stage. Starting
with $M_0=V$, let $M_{l+1}$ be an ordered free amalgam of $M_l$ and $C$ over
$B$, where $(k,B,f)$ is the task met at stage $l\ge k$. Here $B\cle M_k\cle M_l$
by induction and Lemma~\ref{lem:calc}(c), so by Lemma~\ref{lem:calc}(d) the
amalgam is positive, $M_l\cle M_{l+1}$, and the image of $C$ is closed in
$M_{l+1}$. Let $M=\bigcup_l M_l$. By transitivity $M_l\cle M_{l'}$ for
$l\le l'$, hence $M_l\cle M$; this gives (M1). For (M2), choose $k$ with
$B\subseteq M_k$; then $B\cle M_k$, the task
$(k,B,f)$ is met at some stage $l$, and the resulting copy of $C$ is closed in
$M_{l+1}\cle M$, hence in $M$. For (M3) one extends a given isomorphism
$\varphi\colon B\to B'$ of finite closed subsets step by step: for $x\in M$
choose by (M1) a finite closed $C\supseteq B\cup\{x\}$; by (M2), applied to
$B'$ and to $\varphi^{-1}\colon B'\to B\subseteq C$, the map $\varphi$ extends to
an isomorphism of $C$ onto a closed subset of $M$. Alternating between
$\varphi$ and $\varphi^{-1}$ along an enumeration of $M$ gives an automorphism.
\end{proof}

The existence of such limits for classes with a notion of strong
substructure is standard; see
\cite[Theorem~2.6]{EHN}. We fix $M$ as in Proposition~\ref{prop:generic} and put
\[
 G=\Aut(M),\qquad \K=\{B\subseteq M:\ B \text{ finite},\ B\cle M\},
\]
where automorphisms preserve the order and the hyperedges. Positivity gives
$\varnothing\in\K$. We use the notation of Section~\ref{sec:reduction}:
$G_B$ is the pointwise stabilizer of $B$, and $\mathcal A=\mathcal A(M,G,\K)$.

\begin{lemma}\label{lem:AMequiv}
The triple $(M,G,\K)$ satisfies \textup{(K1)}. For $B,C\in\K$ the arrows
$B\to C$ of $\mathcal A$ are exactly the closed embeddings $B\to C$, and
$\mathcal A$ is equivalent to $\Cd$.
\end{lemma}

\begin{proof}
(K1) follows from (M1) and from the invariance of $\cle$ under automorphisms.
If $g\in G$ and $gB\subseteq C$, then $gB\cle M$, so $gB\cle C$ by
Lemma~\ref{lem:calc}(b), and $g|_B$ is a closed embedding. Conversely, if
$f\colon B\to C$ is a closed embedding, then $f(B)\cle C\cle M$, so
$f(B)\in\K$ by Lemma~\ref{lem:calc}(c), and $f\colon B\to f(B)$ extends to an
element of $G$ by (M3). Finally every object of $\Cd$ is isomorphic to a member
of $\K$ by (M2) with $B=\varnothing$, so the functor that sends a non-empty
$B\in\K$ to the object of $\Cd$ isomorphic to it is an equivalence.
\end{proof}

\begin{lemma}\label{lem:K2delta}
$(M,G,\K)$ satisfies \textup{(K2)}: $\Fix_M(G_B)=B$ for every $B\in\K$.
\end{lemma}

\begin{proof}
Let $x\in M\setminus B$ and choose $D\in\K$ containing $B\cup\{x\}$. Let $D'$ be a
second copy of $D$, with $D\cap D'=B$, and form an ordered free amalgam of $D$
and $D'$ over $B$; note that $B\cle D$ by Lemma~\ref{lem:calc}(b). By
Lemma~\ref{lem:calc}(d) the inclusion of $D$ into it is a closed embedding, so
by (M2) the amalgam embeds into $M$ over $D$, with closed image; the image of
$D'$ is then closed in $M$. The isomorphism of $D$ onto this image that is the
identity on $B$ extends by (M3) to some $g\in G_B$, and $gx\ne x$.
\end{proof}

\begin{lemma}\label{lem:SCdelta}
The category $\mathcal A$ satisfies \textup{(SC)}, and $B\mathcal A$ is weakly
contractible.
\end{lemma}

\begin{proof}
Both properties are invariant under equivalence, so we may replace $\mathcal A$
by the following category $\mathcal C$. Fix an infinite set $\Omega$ and
injections $\lambda,\rho\colon\Omega\to\Omega$ with disjoint images. The objects
of $\mathcal C$ are the non-empty positive ordered hypergraphs whose vertex set
is a finite subset of $\Omega$, and its arrows are the closed embeddings; by
Lemma~\ref{lem:AMequiv} it is equivalent to $\mathcal A$.

Let $P$ be an object of $\mathcal C$, let $\mathcal U$ be a non-empty full
cosieve of $P\backslash\mathcal C$, and fix an object $(Z,z\colon P\to Z)$ of
$\mathcal U$. For an object $(X,x\colon P\to X)$ of $P\backslash\mathcal C$ let
$T(X)$ be the hypergraph on $\lambda(X)\cup\rho(Z\setminus z(P))$ that is the free
amalgam of the copy $\lambda(X)$ of $X$ and of the copy
$\lambda(x(P))\cup\rho(Z\setminus z(P))$ of $Z$ over $\lambda(x(P))$, ordered
with $\lambda(X)$ first. The base is closed in both copies because $x$ and
$z$ are closed embeddings, so by Lemma~\ref{lem:calc}(d) $T(X)$ is positive and
the maps
\[
 \eta_X\colon X\to T(X),\ u\mapsto\lambda(u),\qquad
 \gamma_X\colon Z\to T(X),\ z(p)\mapsto\lambda(x(p)),\ w\mapsto\rho(w)\ \
 (w\notin z(P)),
\]
are closed embeddings; they are arrows under $P$. For an arrow
$f\colon X\to X'$ of $P\backslash\mathcal C$ let
$T(f)=\lambda f\lambda^{-1}$ on $\lambda(X)$ and the identity on
$\rho(Z\setminus z(P))$. It is increasing, because $f$ is increasing and
commutes with the maps from $P$, so that it preserves the gap of $P$ in which a
vertex lies. It is an embedding, since a hyperedge of $T(X')$ lies in one of
the two copies and $f$ is an embedding. Its image is closed, because $T(X')$
is the free amalgam of $\lambda(X')$ and of the image of $T(X)$ over
$\lambda(f(X))$, and $f(X)\cle X'$. So $T$ is an endofunctor of
$P\backslash\mathcal C$, and $\eta\colon\id\to T$ and
$\gamma\colon\mathrm{const}_{(Z,z)}\to T$ are natural transformations. Since
$\mathcal U$ is a cosieve containing $(Z,z)$, the arrows $\gamma_X$ show that
$T$ maps $P\backslash\mathcal C$ into $\mathcal U$. The restrictions of $\eta$
and $\gamma$ to $\mathcal U$ connect the identity of $\mathcal U$ to a constant
functor, so $B\mathcal U$ is contractible \cite[\S 1, Proposition~2]{Quillen}.

The same construction with the empty set in the place of $P$ (so that
$T(X)=\lambda(X)\cup\rho(Z)$ is a disjoint union, with $\lambda(X)$ first)
gives, for a fixed object $Z$, an endofunctor $T$ of $\mathcal C$ with natural
transformations $\id\to T\leftarrow\mathrm{const}_Z$. Hence $B\mathcal C$ is
contractible.
\end{proof}

\subsection{Open subgroups}

\begin{lemma}[the bridge]\label{lem:bridgedelta}
Let $D$ and $E$ be positive ordered hypergraphs, let $B\cle D$ and $C\cle D$,
put $F=B\cap C$, and let $d_0,d_1\colon D\to E$ be closed embeddings that agree
on $F$. Then there are a positive ordered hypergraph $Z$, a closed embedding
$e\colon E\to Z$ and closed embeddings $a_0,\dots,a_{2N}\colon D\to Z$ with
\[
 a_0=ed_0,\quad a_{2N}=ed_1,\quad a_{2j}|_B=a_{2j+1}|_B,\quad
 a_{2j+1}|_C=a_{2j+2}|_C\qquad(0\le j<N).
\]
\end{lemma}

\begin{proof}
Put $m=\max(|B|,|C|)$ and take $N$ so large that Lemma~\ref{lem:bridgeset}
applies and
\begin{equation}\label{eq:Nbound}
 s:=2N>3\,\bigl(\delta(E)+2m\bigr),\qquad N\ge 2 .
\end{equation}
Let $Z$, $e$, $a_i$ and the pieces $Z_0,\dots,Z_{s-1}$ be as in
Lemma~\ref{lem:bridgeset}; the indices of pieces are read modulo $s$. Every
piece is a copy of $E$ or of $D$; let the hyperedges of $Z$ be the images of
the hyperedges of the pieces. It remains to prove that every piece is an
induced subhypergraph of $Z$, that $Z$ is positive, and that every piece is
closed in $Z$.

\emph{Pieces are induced.} For two consecutive pieces the intersection
$Z_i\cap Z_{i+1}$ is the common image of $B$ or of $C$ (Lemma~\ref{lem:bridgeset}(iii)), and the two pieces induce
the same hypergraph on it, because $d_0$ and $d_1$ are embeddings. Call this
fact $(\ast)$. Let $h$ be a hyperedge of the piece $Z_j$ whose three vertices lie in
$Z_i$, $i\ne j$. If all three are $F$-points, $h$ is the image of a hyperedge
of $D$ contained in $F$, hence a hyperedge of every piece. Otherwise one of
them lies in at most three cyclically consecutive pieces, among them $Z_i$ and
$Z_j$. If $i$ and $j$ are consecutive, $h$ is a hyperedge of $Z_i$ by $(\ast)$.
If not, then $\{i,j\}=\{1,s-1\}$ by Lemma~\ref{lem:bridgeset}(ii), and every
point of $Z_1\cap Z_{s-1}$ lies in $Z_0$; applying $(\ast)$ to $(Z_j,Z_0)$ and
then to $(Z_0,Z_i)$ shows again that $h$ is a hyperedge of $Z_i$.

\emph{The cut.} Fix an index $r$. Let $L$ be obtained from the disjoint union
of the pieces by making all the identifications of Lemma~\ref{lem:bridgeset}
except those between the pieces $r-1$ and $r$, with hyperedges the images of
the hyperedges of the pieces, and let $q\colon L\to Z$ be the quotient map.
Write $T_0,\dots,T_{s-1}$ for the images in $L$ of the pieces
$r,r+1,\dots,r-1$, in this order. By Lemma~\ref{lem:bridgeset}(iv) a point of
$L$ lies in $T_k$ for a set of consecutive values of $k$. Hence
$T_k\cap(T_0\cup\dots\cup T_{k-1})=T_k\cap T_{k-1}=:A_k$. Applying $(\ast)$ along
the consecutive pieces between two pieces that contain the vertices of a
hyperedge shows, as in the first step, that every $T_k$ is an induced
subhypergraph of $L$. So $L_k=T_0\cup\dots\cup T_k$ is the free amalgam of
$L_{k-1}$ and $T_k$ over $A_k$. The set $A_k$ is a copy of $B$ or of $C$; it is closed in $T_k$ and in
$T_{k-1}$, by hypothesis and, for a piece that is a copy of $E$, by
Lemma~\ref{lem:calc}(c). By induction on $k$ and Lemma~\ref{lem:calc}(c),(d),
$L_k$ is positive and each of $T_0,\dots,T_k$ is closed in $L_k$. So $L$ is
positive and every $T_k$ is closed in $L$. By Lemma~\ref{lem:bridgeset}(iv) the
map $q$ identifies at most $m$ pairs of points, and every hyperedge of $Z$ is
the image of a hyperedge of $L$. Hence, for $U\subseteq Z$ and $\widetilde U=q^{-1}(U)$,
\begin{equation}\label{eq:cutloss}
 \delta_Z(U)\ \ge\ \delta_L(\widetilde U)-2m .
\end{equation}

\emph{Positivity and closedness.} Let $U\subseteq Z$ be non-empty.

\emph{Case 1: some piece $Z_r$ contains no point of $U$ other than
$F$-points.} Cut between $r-1$ and $r$. No point of $U$ is then the image of
two points of $L$, so $q$ maps $\widetilde U$ bijectively onto $U$ and
$E_L(\widetilde U)$ bijectively onto $E_Z(U)$. Hence
$\delta_Z(U)=\delta_L(\widetilde U)>0$. If $U$ properly contains a piece $Z_i$,
then $\widetilde U$ properly contains the corresponding $T_k$, and
$\delta_Z(U)=\delta_L(\widetilde U)>\delta_L(T_k)=\delta_Z(Z_i)$.

\emph{Case 2: every piece contains a point of $U$ that is not an $F$-point.} Cut
between $s-1$ and $0$, and put $\widetilde U_k=\widetilde U\cap T_k$. Applying
\eqref{eq:amalgam} to $L_k=L_{k-1}\cup T_k$ for $k=1,\dots,s-1$ gives
\[
 \delta_L(\widetilde U)=\delta(\widetilde U_0)
 +\sum_{k=1}^{s-1}\bigl(\delta(\widetilde U_k)-\delta(\widetilde U_k\cap A_k)\bigr).
\]
By Lemma~\ref{lem:calc}(b), $\widetilde U_k\cap A_k\cle\widetilde U_k$, so the
$k$-th summand is non-negative, and it is at least $1$ when
$\widetilde U_k\not\subseteq A_k$; also $\delta(\widetilde U_0)\ge 1$ when
$\widetilde U_0\ne\varnothing$. Let $J$ be the set of those $k$ for which
$\widetilde U_k$ contains a point that is not an $F$-point and does not lie in
$T_{k-1}$ (no condition on $T_{k-1}$ for $k=0$). Each $k\in J$ contributes at
least $1$, so $\delta_L(\widetilde U)\ge|J|$. Let $0\le k'<s$. By assumption
$T_{k'}$ contains a point $u\in\widetilde U$ that is not an $F$-point; the set
of $k$ with $u\in T_k$ consists of at most three consecutive values, and its
least element $k$ belongs to $J$, with $k\le k'\le k+2$. Hence $3|J|\ge s$, and
by \eqref{eq:cutloss} and \eqref{eq:Nbound}
\[
 \delta_Z(U)\ \ge\ \frac s3-2m\ >\ \delta(E)\ \ge\ \delta(D)\ >\ 0 ,
\]
where $\delta(E)\ge\delta(D)$ because $d_0(D)\cle E$. In particular
$\delta_Z(U)>\delta_Z(Z_i)$ for every piece $Z_i$.

In both cases $\delta_Z(U)>0$, and $\delta_Z(U)>\delta_Z(Z_i)$ whenever $U$
properly contains the piece $Z_i$. So $Z$ is positive and every piece is closed
in $Z$; with the first step, $e$ and all $a_i$ are closed embeddings.
\end{proof}

\begin{proposition}\label{prop:stabdelta}
For all $B,C\in\K$ we have $\langle G_B,G_C\rangle=G_{B\cap C}$.
\end{proposition}

\begin{proof}
One inclusion is clear. Let $g\in G_{B\cap C}$. Choose $D\in\K$ containing
$B\cup C$ and $E\in\K$ containing $D\cup gD$. Then $B\cle D$ and $C\cle D$ by
Lemma~\ref{lem:calc}(b), and the inclusion $d_0\colon D\to E$ and
$d_1=g|_D\colon D\to E$ are closed embeddings that agree on $B\cap C$. Let
$Z$, $e$ and $a_0,\dots,a_{2N}$ be as in Lemma~\ref{lem:bridgedelta}. By (M2)
there is an embedding $h\colon Z\to M$ with closed image and $he=\id_E$. Each
$ha_i\colon D\to M$ is an isomorphism onto a closed subset, so by (M3)
$ha_i=g_i|_D$ for some $g_i\in G$; since $ha_0=d_0$ and $ha_{2N}=d_1$ we may take
$g_0=1$ and $g_{2N}=g$. The relations between consecutive $a_i$ say that
$g_i^{-1}g_{i+1}$ lies in $G_B$ for $i$ even and in $G_C$ for $i$ odd. Hence
$g=\prod_{i=0}^{2N-1}g_i^{-1}g_{i+1}\in\langle G_B,G_C\rangle$.
\end{proof}

\begin{corollary}\label{cor:opendelta}
$(M,G,\K)$ satisfies \textup{(K3)}: every open subgroup of $G$ is $G_F$ for some
$F\in\K$.
\end{corollary}

\begin{proof}
Let $U$ be an open subgroup. By (K1) it contains $G_B$ for some $B\in\K$;
choose such a $B$ of least cardinality. For $g\in U$ we have
$G_{gB}=gG_Bg^{-1}\le U$, so $G_{B\cap gB}\le U$ by
Proposition~\ref{prop:stabdelta}; here $B\cap gB\in\K$ by
Lemma~\ref{lem:calc}(b),(c). By minimality $gB=B$, and an increasing bijection
of a finite linearly ordered set is the identity, so $g\in G_B$. Hence $U=G_B$.
\end{proof}

\subsection{The two properties}

\begin{theorem}\label{thm:total}
The category $\Cd$ is total: it is non-empty and connected, and for every
functor $L\colon\Cd\to\Set$ every connected component of $\El(L)$ is weakly
contractible.
\end{theorem}

\begin{proof}
By Lemmas~\ref{lem:AMequiv}, \ref{lem:K2delta}, \ref{lem:SCdelta} and
Corollary~\ref{cor:opendelta} the triple $(M,G,\K)$ satisfies (K1)--(K4). By
Proposition~\ref{prop:reduction} every action component of $\mathcal A$ is
weakly contractible or weakly equivalent to $\mathcal A$, and $B\mathcal A$ is
weakly contractible by Lemma~\ref{lem:SCdelta}. Since $\mathcal A$ is equivalent
to $\Cd$, the same holds for $\Cd$, and Proposition~\ref{prop:total} applies.
\end{proof}

\begin{theorem}\label{thm:notexact}
Let $R$ be a non-zero ring. For an object $X$ of $\Cd$ let
$P_X=R[\Cd(X,-)]$, a functor from $\Cd$ to left $R$-modules, and let
\[
 \Theta\colon P_W\longrightarrow P_V\oplus P_V,\qquad
 \Theta(g)=(gv_2-gv_1,\ gv_3-gv_1)\quad(g\colon W\to X).
\]
Then $\Theta$ is a monomorphism, and $\colim\Theta\colon R\to R\oplus R$ is zero.
In particular the functor $\colim$ from $\Cd$-indexed diagrams of $R$-modules to
$R$-modules is not exact; for $R=\Z$, the category $\Cd$ is not exact.
\end{theorem}

\begin{proof}
Let $X$ be an object and let $t=\sum_g c_g\,g$ be a non-zero element of
$P_W(X)$; let $g_1,\dots,g_r$ be the arrows $W\to X$ with $c_g\ne 0$. An
increasing map of $W$ is determined by its image, so the images of the $g_k$
are $r$ distinct hyperedges of $X$. Let $U$ be their union and $v=|U|$. Then
$0<\delta(U)\le 2v-3r$. The $r$ hyperedges have $3r<2v$ incidences with the $v$
points of $U$, each of which lies on at least one of them; so some point $u$
lies on exactly one, say on the image of $g_k$, where $u=g_k(i)$. The arrow
$g_kv_i\colon V\to X$ has image $\{u\}$, and no other arrow $g_lv_j$ with
$c_{g_l}\ne 0$ has this image. The coefficient of $g_kv_i$ is therefore
$c_{g_k}$ in the first component of $\Theta(t)$ if $i=2$, it is $c_{g_k}$ in the
second component if $i=3$, and it is $-c_{g_k}$ in both if $i=1$. So
$\Theta(t)\ne 0$.

The colimit of a representable functor $P_X$ is $R$, the canonical map
$P_X(Y)\to R$ being the augmentation $\sum c_f f\mapsto\sum c_f$. Both
components of $\Theta(g)$ have augmentation zero, so $\colim\Theta=0$, which is
not injective because $\colim P_W=R\ne 0$.
\end{proof}
